\chapter{2010年浙江卷（文）}

\section{单选题}

\begin{problem}
设 \(P = \{x \mid x < 1\}\)，\(Q = \{x \mid x^2 < 4\}\)，则 \(P \cap Q =\)（\quad）

\choices
    {\(\{x\mid -1 < x < 2\}\)}
    {\(\{x\mid -3 < x < - 1\}\)}
    {\(\{x\mid 1 < x < 4\}\)}
    {\(\{x\mid -2 < x < 1\}\)}
\end{problem}

\begin{problem}
已知函数 \( f(x) = \log_6(x + 1) \)，若 \( f(\mathrm{o}) = 1 \)，则 \( \alpha = \)（\quad）

\choices
    {0}
    {1}
    {2}
    {3}
\end{problem}

\begin{problem}
设 \(1\) 为虚数单位，则 \(\frac{5 - \i}{1 + \i} =\)（\quad）

\choices
    {\(-2\) -3i}
    {\(-2 + 3\i\)}
    {2 - 3i}
    {\(2 + 3\i\)}
\end{problem}

\begin{problem}
某程序框图如图所示，若输出的 S = 57，则判断框内为（\quad）

\choices
    {\(k > 4\)}
    {\(k > 57\)}
    {\(k > 67\)}
    {\(k > 77\)}
\end{problem}

\begin{problem}
设 \(S_{n}\) 为等比数列 \(\{a_{n}\}\) 的前 \(n\) 项. 知 \(k_{2} + a_{5} = 0\)，则 \(\frac{S_{5}}{S_{2k}} =\)（\quad）

\choices
    {\(-11\)}
    {\(-8\)}
    {5}
    {11}
\end{problem}

\begin{problem}
设 \(0 < x < \frac{\pi}{2}\)，则“\(x \sin^2 x < 1\)”是“\(x \sin x < 1\)”的（\quad）

\choices
    {充分而不必要条件}
    {必要而不充分条件}
    {充分必要条件}
    {既不充分也不必要条件}
\end{problem}

\begin{problem}
若实数 \(x, y\) 满足不等式组 \(\left\{ \begin{array}{l} x + 3y - 3 \geqslant 0, \\ 2x - y - 3 \leqslant 0, \\ x - y + 1 \geqslant 0, \end{array} \right.\) 则 \(x + y\) 的最大值为（\quad）

\choices
    {9}
    {\(\frac{15}{7}\)}
    {1}
    {\(\frac{7}{15}\)}
\end{problem}

\begin{problem}
若某几何体的三视图 (单位: cm) 如图所示，则此几何体的体积是（\quad）

正视图

侧视图

俯视图（\quad）

\choices
    {\(\frac{352}{3}\mathrm{cm}^3\)}
    {\(\frac{320}{3}\mathrm{cm}^3\)}
    {\(\frac{224}{3}\mathrm{cm}^3\)}
    {\(\frac{160}{3}\mathrm{cm}^3\)}
\end{problem}

\begin{problem}
已知 \(x_0\) 是函数 \(f(x) = 2^x + \frac{1}{1 - x}\) 的一个零点. 若 \(x_1 \in (1, x_0)\)，\(x_2 \in (x_0, +\infty)\)，则（\quad）

\choices
    {\(f(x_{1}) < 0, f(x_{2}) < 0\)}
    {\(f(x_{1}) < 0, f(x_{2}) > 0\)}
    {\(f(x_{1}) > 0, f(x_{2}) < 0\)}
    {\(f(x_{1}) > 0, f(x_{2}) > 0\)}
\end{problem}

\begin{problem}
设 O 为坐标原点，\( F_{1} \) ，\( F_{2} \) 是双曲线 \( \frac{x^{2}}{a^{2}} - \frac{y^{2}}{b^{2}} = 1 (a > 0, b > 0) \) 的焦点. 若在双曲线上存在点 P，满足 \( \angle F_{1}PF_{2} = 60^{\circ} \) ，\( |OP| = \sqrt{7}a \) ，则该双曲线的渐近线方程为（\quad）

\choices
    {\(x \pm \sqrt{3} y = 0\)}
    {\(\sqrt{3} x \pm y = 0\)}
    {\(x \pm \sqrt{2} y = 0\)}
    {\(\sqrt{2} x \pm y = 0\)}
\end{problem}

\section{填空题}

\begin{problem}
在如图所示的茎叶图中，甲、乙两组数据的中位数分别是，. 1 2 3 8 2 9 6 9 1 3 5 2 5 4 8 7 8 5 5 6 6 7 （12）函数 \( f(x) = \sin^2\left(2x - \frac{\pi}{4}\right) \) 的最小正周期是 \fillinblank{} \fillinblank{}.
\end{problem}

\begin{problem}
已知平面向量 \(\alpha, \beta\) 满足 \(|\alpha| = 1, |\beta| = 2, \alpha \perp (\alpha - 2\beta)\)，则 \(|2\alpha + \beta|\) 的值是 \fillinblank{} \fillinblank{}.
\end{problem}

\begin{problem}
在如下数表中，已知每行、每列中的数都成等差数列. 第1列 第2列 第3列 ... 第1行 1 2 3 ... 第2行 2 4 6 ... 第3行 3 6 9 ... ... ... ... ... ... 那么，位于表中的第 n 行第 \( n+1 \) 列的值是 \fillinblank{} \fillinblank{}.
\end{problem}

\begin{problem}
若正实数 \(x, y\) 满足 \(2x + y + 6 = xy\)，则 \(xy\) 的最小值是 \fillinblank{} \fillinblank{}.
\end{problem}

\begin{problem}
某商家一月份至五月份累计销售额达 3860 万元，预测六月份销售额为 500 万元. 七月份销售额比六月份递增 2\%，八月份销售额比七月份递增 \(-2\)\%，九、十月份销售总额与七、八月份销售总额相等. 若一月至十月份销售总额至少达 7000 万元，则 x 的最小值是 \fillinblank{} \fillinblank{}.
\end{problem}

\begin{problem}
在平行四边形 \(ABCD\) 中，\(O\) 是 \(AC\) 与 \(BD\) 的交点，\(P, Q, M, N\) 分别是线段 \(OA, OB, OC, OD\) 的中点. 在 \(A, P, M, C\) 中任取一点记为 \(E\)，在 \(B, Q, N, D\) 中任取一点记为 \(F\). 设 \(G\) 为满足向量 \(\overrightarrow{OG} = \overrightarrow{OE} + \overrightarrow{OP}\) 的点，则在上述的点 \(G\) 组成的集合中的点，落在平行四边形 \(ABCD\) 外 (不含边界) 的概率为 \fillinblank{} \fillinblank{}.
\end{problem}

\section{解答题}

\begin{problem}
在 \(\triangle ABC\) 中，角 \(A, B, C\) 所对的边分别为 \(a, b, c\). 设 \(S\) 为 \(\triangle ABC\) 的面积，满足 \(S = \frac{\sqrt{3}}{2} (a^2 + b^2 - c^2)\) (1) 求角 C 的大小; (2) 求 \( \sin A + \sin B \) 的最大值.
\end{problem}

\begin{problem}
设 \(a_1, d\) 为实数. 首项为 \(a_1\)，公差为 \(d\) 的等差数列 \(\{a_n\}\) 的前 \(n\) 项和为 \(S_n\)，满足 \(S_5 S_6 + 15 = 0\). (1) 若 \( S_{5}=5 \) ，求 \( S_{6} \) 及 \( a_{1} \) ; (2) 求 d 的取值范围.
\end{problem}

\begin{problem}
已知函数 \( f(x) = (x - a)^2 (x - b) \) (\( a, b \in \R, a < b \)). (1) 当 \(a = 1, b = 2\) 时，求曲线 \(y = f(x)\) 在点 \((2, f(2))\) 处的切线方程； (2) 设 \(x_{1}, x_{2}\) 是 \(f(x)\) 的两个极值点，\(x_{3}\) 是 \(f(x)\) 的一个零点，且 \(x_{1} \neq x_{1}\)，\(x_{3} \neq x_{2}\). 证明: 存在实数 \(x_{4}\)，使得 \(x_{1}, x_{2}, x_{3}, x_{4}\) 按某种顺序排列后成等差数列，并求 \(x_{4}\).
\end{problem}

\begin{problem}
已知 \(m\) 是非零实数，抛物线 \(C: y^{2} = 2px (p > 0)\) 的焦点 \(F\) 在直线 \(l: x - my - \frac{m^{2}}{2} = 0\) 上. (1) 若 \(m = \frac{2}{2}\)，求抛物线 \(C\) 的方程; (2) 设直线 \(l\) 与抛物线 \(C\) 交于 \(A, B\) 两点，过 \(A, B\) 分别作抛物线 \(C\) 的准线的垂线，垂足为 \(A_{1}, B_{1}, \triangle A_{1}F, \triangle BB_{1}F\) 的重心分别为 \(G, H\). 求证：对任意非零实数 \(m\)，抛物线 \(C\) 的准线与 \(x\) 轴的交点在以线段 \(GH\) 为直径的圆外.
\end{problem}

\begin{problem}
如图，在平行四边形 ABCD 中， \( AB = 2BC, \angle ABC = 120^{\circ} \) . E 为线段 AB 的中点，将 \( \triangle ADE \) 沿直线 DE 翻折成 \( \triangle A'DE \) ，使平面 \( A'DE \perp \) 平面 BCD，F 为线段 \( A'C' \) 的中点. (1) 求证: \( BF // \) 平面 \( A'DE \) ; (2) 设 M 为线段 DE 的中点，求直线 FM 与平面 \( A^{\prime}DE \) 所成角的余弦值.
\end{problem}

